\chapter{Introducción}
\label{cap:introduccion}

% \chapterquote{La revolución industrial y sus consecuencias han sido un desastre para la raza humana}{Theodore Kaczynski}

% En los últimos años el avance de tecnologías inmersivas como la \gls{vr} ha proporcionado un gran cambio en la forma en la que los usuarios interactúan en los entornos digitales.
% Grandes empresas como Meta o Apple han apostado por el desarrollo de hardware para estas tecnologías con el lanzamiento de Meta Quest 3 y Apple Vision Pro, así como con la creación de plataformas de interacción social en el \gls{metaverso} con el desarrollo de Meta Horizon Worlds.
% Esta tecnología ha abierto nuevas oportunidades para explorar modos de comunicación más naturales en espacios virtuales mediante la comunicación no verbal.

% La comunicación no verbal es aquella que no se basa en el uso de palabras, sino que utiliza las imágenes y los gestos para transmitir un mensaje.
% En nuestro día a día representa una parte esencial de las comunicaciones humanas, sin embargo su representación en entornos virtuales es más limitada y dependiente de acciones predefinidas como animaciones o ``emotes'' ya introducidos en el juego.
% Es por esto que el desarrollo de herramientas que permitan capturar e interpretar en tiempo real estos gestos constituye un gran paso hacia entornos virtuales más expresivos y realistas.

Los videojuegos están compuestos por múltiples elementos que interactúan entre sí para construir una experiencia coherente e inmersiva para el jugador, entre los que se encuentran el avatar del jugador, los escenarios, los objetos interactivos y, de forma habitual, los personajes no jugables.
Estos personajes no jugables (NPC, por sus siglas en inglés de Non-Player Character) desempeñan un papel relevante en la construcción de mundos virtuales, ya que son responsables de proporcionar interacciones, responder a las acciones del jugador y desarrollar comportamientos coherentes con el contexto del juego.

El desarrollo de estos comportamientos constituye una tarea que involucra tanto aspectos de diseño como de implementación.
Habitualmente, los diseñadores determinan las acciones, rutinas e interacciones que debe desarrollar un NPC, mientras que los programadores se encargan de trasladar estas especificaciones a una implementación ejecutable.
Para ello pueden emplearse diferentes técnicas de inteligencia artificial, entre las que destacan las máquinas de estados finitos (Finite State Machines, FSM) y los árboles de comportamiento (Behavior Trees, BT), siendo estos últimos el principal objeto de estudio de este trabajo.

El desarrollo de comportamientos para NPCs es, además, un proceso inherentemente iterativo.
Durante la implementación es necesario comprobar el comportamiento obtenido, detectar situaciones no previstas y modificar la lógica para adaptarla a los requisitos del juego.
Como consecuencia, diseñadores y programadores pueden verse obligados a repetir sucesivamente las fases de diseño, implementación, evaluación y modificación de los comportamientos.
En el caso de los árboles de comportamiento, estos cambios implican con frecuencia crear, reorganizar o modificar nodos, condiciones y relaciones de ejecución, lo que puede incrementar el esfuerzo necesario a medida que aumenta la complejidad del comportamiento.

A este coste de desarrollo se añade la necesidad de disponer de conocimientos técnicos específicos para construir y modificar correctamente estas estructuras.
La utilización de un sistema basado en árboles de comportamiento requiere comprender tanto su lógica de ejecución como los elementos que lo componen y sus relaciones, lo que puede dificultar su utilización por parte de desarrolladores con menor experiencia en inteligencia artificial o por perfiles no especializados en programación.

En este contexto, la generación automática de árboles de comportamiento a partir de descripciones realizadas en lenguaje natural constituye una posible vía para reducir el esfuerzo asociado a su creación y modificación.
Permitir que un usuario pueda describir el comportamiento deseado mediante una especificación en lenguaje natural y obtener de forma automática una estructura ejecutable podría reducir parte del trabajo manual necesario durante el desarrollo, facilitar la iteración sobre los comportamientos y disminuir las barreras técnicas asociadas a su implementación.

La aparición de los modelos grandes de lenguaje (Large Language Models, LLMs), capaces de interpretar instrucciones en lenguaje natural y generar estructuras complejas a partir de ellas, abre una nueva posibilidad para automatizar este proceso.
Por este motivo, este trabajo estudia la utilización de modelos de lenguaje para la generación automática de árboles de comportamiento destinados al control de personajes no jugables, analizando las posibilidades y limitaciones de este enfoque en el contexto del desarrollo de videojuegos.

\section{Motivación}
\label{sec:Motivacion}
% Como se ha mencionado en la introducción, el trabajo de los diseñadores a la hora de implementar BTs es un proceso lento y tedioso que se puede repetir varias veces debido a rediseños de los personajes.
% Además, con los avances tecnológicos con consolas y equipos de mayor potencia (tanto de gráfica como de memoria), los videojuegos se están volviendo cada vez más grandes y exigentes, aumentando de esta forma la complejidad de los comportamientos de los personajes repercutiendo en el aumento de la complejidad de los \glspl{BT}.
% Debido a estos motivos la generación de los árboles de comportamiento es una línea investigativa que, como se va a expandir en el apartado \ref{sec:GenCom}, se lleva investigando desde hace varios años.
% Además, los avances en los grandes modelos de lenguaje permiten avances en este ámbito, permitiendo el uso del lenguaje natural para la generación  y acercando de esta forma el diseño a usuarios no técnicos.

% Es por los motivos anteriores que este trabajo se inscribe dentro de esta línea de investigación, buscando aportar una herramienta funcional y sencilla de usar para todo tipo de usuarios.

% Para ello se describen a continuación los diferentes obejtivos que se buscan en este trabajo.

La creciente complejidad de los videojuegos ha incrementado también la complejidad de los comportamientos que deben implementar sus personajes no jugables.
A medida que se incorporan un mayor número de acciones, situaciones y posibilidades de interacción, las estructuras encargadas de controlar estos comportamientos requieren un mayor número de elementos y relaciones entre ellos.
En este contexto, los árboles de comportamiento presentan características como la modularidad y la escalabilidad que resultan especialmente adecuadas para representar comportamientos complejos (\cite{DEFCON}).

Sin embargo, estas características no eliminan el coste asociado a la creación y mantenimiento de los árboles.
La construcción manual de comportamientos puede resultar laboriosa y requerir conocimientos especializados, especialmente cuando los comportamientos deben modificarse o ampliarse de manera iterativa.
La literatura señala precisamente el esfuerzo que supone para los diseñadores desarrollar comportamientos de NPC y proponen mecanismos destinados a automatizar su expansión (\cite{CBR}).
Asimismo, hay artículos que ponen de manifiesto la necesidad de facilitar la interacción entre los aspectos de diseño y los aspectos técnicos de la implementación de comportamientos (\cite{BTValRef}).

La dificultad de autoría de los árboles de comportamiento también ha motivado la investigación de herramientas destinadas a facilitar su creación.
Hay autores que proponen una herramienta de autoría basada en una interfaz visual que busca simplificar la construcción de estos árboles y acercar su edición al proceso de diseño (\cite{AIPaint}).
Estos trabajos muestran que la reducción del esfuerzo necesario para especificar y modificar comportamientos constituye una problemática anterior a la aparición de los LLMs.

Como consecuencia, se han desarrollado diferentes aproximaciones para automatizar total o parcialmente la construcción de los árboles de comportamiento.
Entre ellas se encuentran métodos basados en la reutilización de comportamientos, técnicas de evolución automática y métodos de síntesis a partir de especificaciones formales, demostrando que la generación de árboles de comportamiento constituye una línea de investigación consolidada.
Estas aproximaciones permiten reducir parte del trabajo manual, aunque generalmente requieren representaciones estructuradas del problema, funciones de evaluación, demostraciones o conocimientos específicos del dominio.

La aparición de los LLMs introduce una nueva posibilidad para abordar este problema, ya que estos modelos permiten procesar instrucciones expresadas en lenguaje natural y generar estructuras a partir de ellas.
En lugar de requerir que el usuario defina directamente la estructura interna del árboles de comportamiento, puede plantearse que describa el comportamiento deseado mediante lenguaje natural y que el sistema se encargue de realizar la traducción a una representación estructurada.
Esta aproximación permite explorar una interacción más cercana al proceso de diseño y puede reducir las barreras técnicas asociadas a la autoría de comportamientos.

En este contexto, este trabajo se sitúa en la intersección entre la investigación sobre generación automática de los árboles de comportamiento y los avances recientes en LLMs.

\section{Objetivos}
\label{sec:Objetivos}

El objetivo general de este proyecto es la implementación de una herramienta que puedan usar de forma sencilla e intuitiva diseñadores, tanto expertos como los que no tienen una base técnica, para poder generar árboles de comportamiento de una manera más eficaz que los que puedan generar ellos y con una calidad similar a la humana con el fin de aliviar su trabajo.
No se busca sustituir el trabajo de un diseñador, sino aportarles una base útil que les devuelva unas estructuras fácilmente modificables para matizar ciertos aspectos o cambiar de forma sencilla los comportamientos generados.
Para cumplir estos objetivos se han propuesto las siguientes metas durante el trabajo:

\begin{itemize}
	\item Implementación de una pila de tareas y condicionales necesarias para generar diferentes comportamientos.
	\item Diseño e implementación de diferentes escenarios en los que se pueda probar en tiempo real y de forma visual los comportamientos buscados.
	\item Comparar diferentes métodos de implementación y modelos de lenguaje (tiempos de ejecución y calidad de las generaciones) para elegir el idóneo para cumplir el objetivo general.
	\item Implementación de la herramienta generadora usando diferentes técnicas de IA que consigan una generación óptima.
	\item Implementación un sistema de validación de los árboles generados para comprobar que son válidos en cuanto a respeto de dependencias de nodos y cumplimiento de objetivos.
	\item Realizar un plan de pruebas para validar el cumplimiento de los objetivos generales.
\end{itemize}

\section{Plan de trabajo}
El plan de trabajo consiste en varios pasos:
\begin{enumerate}
	%\renewcommand{\theenumi}{\alph{enumi}}
	\item \texttt{Noviembre - Febrero}: Desarrollo e implementación de varios escenarios en un motor de videojuegos que contengan lo necesario para la generación: una pila de tareas y condicionales, capacidad para ejecutar los BTs generados y una forma de comunicarse con la herramienta de generación.
	\item \texttt{Febrero - Abril}: Implementación de una herramienta que, mediante grandes modelos de lenguaje, sea capaz de generar BTs a través del lenguaje natural capaces de satisfacer un comportamiento buscado.
	\item \texttt{Abril - Junio}: Desarrollo de un sistema de validación automática de los BTs generados.
	\item \texttt{Julio - Agosto}: Diseño y posterior ejecución de pruebas con usuarios objetivo para medir el correcto cumplimiento de los objetivos.
\end{enumerate}

\section{Metodología y tecnología}
Para desarrollar este proyecto se han usado diferentes herramientas comúnmente utilizadas en el desarrollo de videojuegos con el objetivo público en mente.

\begin{itemize}
	\item \texttt{Unreal Engine 5.7}: Unreal Engine es uno de los motores líderes de la industria, caracterizado por su robustez técnica.
	Además incorpora un sistema nativo de árboles de comportamiento dirigidos por eventos y un componente de Blackboard, necesario para la interacción del personaje no jugable con el entorno.
	Este motor ha sido seleccionado por su amplio uso en la industria, facilitando que los diseñadores hacia los que está enfocada la herramienta se encuenten en un entorno conocido.
	En esta herramienta se ha creado una pila tareas y condicionales necesarias para la creación de comportamientos, así como varios escenarios para probar las generaciones de los árboles de comportamiento.
	\item \texttt{Python 3}: Este es un lenguaje de programación de alto nivel ampliamente extendido debido a su versatilidad y posición como uso común en el desarrollo de inteligencia artificial.
	Al ser el backend del proyecto, y siguiendo con los objetivos del uso de esta herramienta, la sección desarrollada con este lenguaje es completamente opaca al usuario.
	La herramienta desarrollada con Python es quien se encarga de toda la generación de los árboles, gracias a módulos disponibles de orquestación de LLMs, siendo Langchain el que ha sido usado en este trabajo .
	\item \texttt{Hugging Face}: Es un ecosistema de código abierto para la distribución de modelos de Machine Learning, librerias y datasets del Procesamiento de Lenguaje Natural.
	En este trabajo ha sido usado para crear un sistema RAG, siendo el origen del repositorio de donde se han sacado los datos y quien provee el modelo de embeddings.
	\item \texttt{Mistral AI}: Es una compañía de modelos de inteligencia artificial generativa que ofrece modelos grandes de lenguaje.
	Uno de los modelos usados a lo largo del trabajo ha sido Mistral Large 3 a través de la API de esta empresa.
	Este modelo cuenta con 675B de parámetros totales y 41B activos, y está diseñado como un sistema de ``Mezcla de Expertos'', haciéndolo un modelo robusto que puede cumplir con el propósito de la herramienta.
	\item \texttt{Ollama}: Se trata de una herramienta de código abierto diseñada para gestionar y ejecutar modelos de lenguaje de manera local, garantizando una baja latencia y la ausencia de costes por llamadas de APIs.
	Este entorno hospeda y ejecuta localmente otro de los modelos usados en el trabajo, el modelo llama3:8B, el cual cuenta con 8B de parámetros.
	Al ser un modelo más ligero no es tan preciso como el mencionado previamente, pero su uso ha resultado ser efectivo en este trabajo.
\end{itemize}

\section{Estructura del documento}
Este documento ha sido estructurado de acuerdo con el plan de trabajo mencionado previamente:

\begin{enumerate}
	\item \texttt{Capítulo 2: Estado de la cuestión}. En este capítulo se presenta el marco teórico y tecnológico en el que se encuentra este trabajo.
	En este capítulo se repasan diferentes métodos de generación de comportamientos usadas en videojuegos, tanto su definición como su evolución en la industria y academia, remarcando el interés especial en los árboles de comportamiento.
	También se cuenta la evolución y el contexto actual de los modelos de lenguaje, para, finalmente, contextualizar sobre la línea de investigación del proyecto de la generación automática de comportamientos.
	\item \texttt{Capítulo 3: Herramienta de generación de árboles de comportamiento.} En este capítulo se describe el desarrollo de la herramienta implementada en este trabajo.
	Este capítulo se divide en cuatro secciones:
	\begin{enumerate}
		\item \texttt{Diseño del pipeline del modelo propuesto para la generación de árboles de comportamiento}: Se describe el diseño del flujo de la herramienta ideado para satisfacer los objetivos propuestos.
		\item \texttt{Descripción del entorno en un motor de videojuegos}: Se detalla el entorno creado en un motor de videojuegos para la implementación de varias tareas y condicionales necesarios para generar diferentes tipos de comportamiento, así como el desarrollo de varios escenarios donde se pueden probar los árboles de comportamiento.
		\item \texttt{Herramienta de generación de árboles de comportamiento}: Se explica la implementación de la herramienta que genera un archivo JSON con un árbol de comportamiento desde la descripción del comportamiento buscado a través del lenguaje natural.
		\item \texttt{Validación automática de los árboles de comportamiento}: Se describen los métodos diseñados para la validación de los árboles de comportamiento: la validación del cumplimiento de dependencias entre tareas y del cumplimiento de un objetivo. Asimismo se describen varias pruebas que se han hecho para comprobar la validez del sistema.
	\end{enumerate}
	\item \texttt{Capítulo 4: Evaluación y resultados}. En este capítulo se relatan las pruebas de evaluación con usuarios diseñadas y ejecutadas para validar una serie de hipótesis, así como se repasan los resultados dados.
	\item \texttt{Capítulo 5: Discusión}. En este capítulo se discute la implementación y los resultados dados por todas las secciones del proyecto, así como se comentan las diferentes limitaciones encontradas durante su desarrollo.
	\item \texttt{Capítulo 6: Conclusiones}. En este capítulo se detallan las conclusiones que han salido de los resultados del proyecto, así como se plantea el trabajo a futuro dadas las limitaciones.
\end{enumerate}

\section{Enlace al repositorio}
El código fuente de este trabajo se encuentra en un repositorio público de Github.
La estructura del repositorio es de la siguiente manera:

\begin{itemize}
	\item \texttt{Análisis}: Se encuentra el CSV de las respuestas de las pruebas de usuario y el \textit{Jupyter Notebook} con el código para generar las gráficas (ver sección \ref{sec:Resultados}).
	\item \texttt{bt-dataset-parser}: Proyecto de Python en donde se descarga y se procesa el dataset necesario para mostrar el sistema RAG (ver sección \ref{sec:GeneracionBTs}).
	\item \texttt{BTGenerator}: Proyecto de Unreal que ha servido como base para los entornos (ver sección \ref{sec:DescripcionEntorno}).
	\item \texttt{llm-bt-generator}: Proyecto de Python con la herramienta generadora de BTs (ver sección \ref{sec:GeneracionBTs}).
	\item \texttt{Memoria}: Cófigo fuente en LaTeX de este documento.
\end{itemize}

El enlace al repositorio es el siguiente: \url{https://github.com/pablosm196/TFM}.

\section{Uso de Inteligencia Artificial}
Este documento ha sido escrito con el apoyo puntual de herramientas de Inteligencia Artificial generativa en tareas de redacción.
Su único uso ha consistido en cimentar las ideas establecidas originalmente en el documento reformulando ciertos párrafos, que posteriormente han sido revisados y editados.
En todo caso la estructura del documento, así como las ideas que transimte y la responsabilidad final sobre el contenido recae en el autor.